\documentclass[12pt]{article}
\usepackage[utf8]{inputenc}
\usepackage[T1]{fontenc}
\usepackage{hyperref}
\title{South Wemale Language Resources for K-12 Curriculum Development}
\author{Compiled by Claude Sonnet 5 (claude-sonnet-5) for Andr\'as Kornai}
\date{}
\begin{document}
\maketitle

\section{Introduction}
South Wemale (ISO 639-3 tlw; Glottocode sout2892) is spoken on Seram Island, Maluku, Indonesia. The identification line in the task gives the family as ``Bookkeeping'', which is not a language family and looks like a data error. The sources found describe Wemale as Austronesian (Central Maluku). The search results treat Wemale as a language with northern and southern varieties. Estimates for the southern variety were about 3,700 speakers, but sources disagree: Joshua Project, citing Ethnologue, gives 11,000 for the South Wemale people group. Some sources rate vitality as vigorous, while others suggest the language is moving toward endangerment. Little educational material was found.

\section{Primary Educational Resources}
No primers, textbooks or school programmes for South Wemale were found.

\section{Linguistic and Reference Materials}
\begin{itemize}
\item Wikipedia, ``Wemale language'': an overview. It covers the whole Wemale language, not only the southern variety. \url{https://en.wikipedia.org/wiki/Wemale_language}
\item Wikipedia, ``Wemale people'': ethnographic background. \url{https://en.wikipedia.org/wiki/Wemale_people}
\end{itemize}
The search summary also cites Collins (1983), \emph{The Historical Relationships of the Languages of Central Maluku}, as a comparative reference. No URL for it was found.

\section{Religious and Community Resources}
\begin{itemize}
\item Joshua Project, ``South Wemale in Indonesia'': a people-group profile that reports the group as mostly Christian. \url{https://joshuaproject.net/people_groups/15908}
\item Bakat (sabda.org) database entry for Wemale (weo): a language record with a vitality rating. Its exact contents were not verified beyond the search listing. \url{https://bakat.sabda.org/697-weo-001.html}
\end{itemize}
No scripture translation was confirmed.

\section{Digital and Audio Resources}
No audio, video, corpus or software resources were found.

\section{Curriculum Development Context}
Very little is available. Most speakers also use Indonesian, which is the language of school and trade. A native-language curriculum would need new documentation. The first step is to resolve which variety is meant, since sources mix Wemale as a whole with its northern and southern varieties.
\end{document}